% !TEX root = userguide.tex

\def\headdate{19th February 2015}

\section{Viscoelastic fluid flow using SUPG and omitting DEVSS}
\label{sec:nodevss}

In this section some examples using the equations for viscoelastic fluid flow
will be given after the theory in the next section.
All examples use streamline-upwind Petrov-Galerkin (SUPG), but DEVSS has been
omitted.

\subsection{Basic formulation for G/SUPG}

Recall the explicit stress formulation for the momentum and mass balance of Sec.~\ref{sec:explicitstress}:
\begin{align}
      -\nabla\cdot\big(2\eta_s\ten D(\vek u^{n+1})\big) +\nabla p^{n+1} 
  \blue{-\alpha\nabla\cdot\big(\nabla\vek u^{n+1} -\ten G^T{}^{n+1} \big)}
&=\nabla\cdot\ten \tau(\ten c^{n+1})
        + \rho\vek b \\
     \nabla\cdot\vek u^{n+1}&=0\\
   -\nabla\vek u^{n+1} + \ten G^T{}^{n+1} &= \ten 0
\end{align}
where $\ten c^{n+1}$ has already been computed from a time step update of
the constitutive equation,
and the implicit stress formulation of Sec.~\ref{sec:implicitstress}:
\begin{align}
    -\nabla\cdot\big(2\eta_s\ten D(\vek u^{n+1})\big) +  \nabla p^{n+1} 
  \blue{-\alpha\nabla\cdot\big(\nabla}&\blue{\vek u^{n+1} -\ten G^T{}^{n+1} \big)} \notag\\
  -\nabla\cdot\big(G\Delta t(-\vek u^{n+1}&\cdot\nabla\ten c^n+
        (\nabla\vek u)^T{}^{n+1}\cdot\ten c^n+
      \ten c^n\cdot\nabla\vek u^{n+1})\big)\notag\\
&=\nabla\cdot( G\ten h(\ten c^n,\Delta t) -G\ten 1)
        + \rho\vek b\\
     \nabla\cdot\vek u^{n+1}&=0\\
   -\nabla\vek u^{n+1} + \ten G^T{}^{n+1} &= \ten 0
\end{align}
which both are (generalized) Stokes problems for $(\ten G^{n+1},\vek u^{n+1},p^{n+1})$. 
The DEVSS terms, highlighted in \blue{blue}, stabilize the velocity/pressure/stress space of the
mixed Stokes system  (LBB stable) \cite{Guenette95}. 

The DEVSS stabilization is required for a pure mixed problem,
i.e.\ $\eta_s=0$ and no implicit stress terms (as in the implicit stress formulation above). However, if $\eta_s\neq0$
and/or the stress implicit terms are present, the velocity/pressure/stress space is regularized and the LBB-stabilization by DEVSS is, strictly speaking, not required. In \cite{Baltussen2010} the implicit stress formulation has been used with the DEVSS terms omitted and it seems to work just fine. Note, that in this way the momentum  and mass balance become a $(\vek u^{n+1},p^{n+1})$ Stokes type system. Since, $\ten G^{n+1}$ is still required in updating the constitutive model (G-method), it can be obtained easily after the $(\vek u^{n+1},p^{n+1})$-solve by 
a weak formulation (projection) of 
\[
   \ten G^{n+1} = (\nabla\vek u^{n+1})^T
\]
Remarks:
\begin{enumerate}
   \item The removal of $\ten G^{n+1}$ from the Stokes system can lead to substantial savings in CPU time, especially in 3D.
   \item The Stokes system becomes better conditioned and convergence for iterative methods is probably more easily achieved.
   \item Experience so far, with the omission of DEVSS is limited. For the simple Couette flow (see Sec.~\ref{sec:couette_nodevss})  the stability for 
   the implicit stress formulation
   with $\eta_s=0$ is reduced, i.e.\ the required time step for stability is smaller and also reduces with mesh refinement. For explicit stress and non-zero solvent $\eta_s\neq0$, however, the stability is less affected and the behavior of the perturbation even seems to be improved. 
   \item It is advised to test your problem both with and without DEVSS to decide which works best for you.
\end{enumerate}

\subsection{Planar Couette flow for a Weissenberg number (Wi) of 10: no DEVSS}
\label{sec:couette_nodevss}

(These examples are based on an example developed by Gaetano d'Avino of 
the University of Naples)

We consider an Oldroyd-B (\texttt{couette10.f90}, \texttt{couette12.f90}) and a UCM fluid (\texttt{couette11.f90},
\texttt{couette13.f90})
for a planar Couette problem on a unit square.
The problem is similar to the problem for an Oldroyd-B fluid in
Sec.~\ref{sec:couette} (explicit stress formulation) and a UCM model in Sec.~\ref{sec:couette_IS} (implicit stress formulation) but now omitting DEVSS. Time integration is first-order accurate (\texttt{couette12.f90},
\texttt{couette13.f90}) and second-order accurate (\texttt{couette10.f90},
\texttt{couette11.f90}).



